% Standalone problem document, generated from the adjacent README.md.
% Compile with XeLaTeX. Mathematical content is maintained in Markdown.
\documentclass[11pt,a4paper]{article}
\usepackage[fontsize=10.5pt]{scrextend}
\usepackage[margin=22mm,headheight=15pt,headsep=9mm,footskip=11mm]{geometry}
\usepackage{fontspec}
\setmainfont{texgyrepagella}[Extension=.otf,UprightFont=*-regular,BoldFont=*-bold,ItalicFont=*-italic,BoldItalicFont=*-bolditalic]
\setsansfont{texgyreheros}[Extension=.otf,UprightFont=*-regular,BoldFont=*-bold,ItalicFont=*-italic,BoldItalicFont=*-bolditalic]
\setmonofont{texgyrecursor}[Extension=.otf,UprightFont=*-regular,BoldFont=*-bold,ItalicFont=*-italic,BoldItalicFont=*-bolditalic,Scale=MatchLowercase]
\usepackage{amsmath,amssymb}
\usepackage{unicode-math}
\setmathfont{texgyrepagella-math.otf}
\usepackage{xcolor,microtype,enumitem,needspace,fancyhdr}
\usepackage{bookmark}
\definecolor{ink}{HTML}{17344B}
\definecolor{accent}{HTML}{197D80}
\definecolor{muted}{HTML}{596773}
\hypersetup{colorlinks=true,linkcolor=accent,urlcolor=accent,pdftitle={AV-02 - Hardness
of the spectral-norm condition
number},pdfauthor={Open Problems in Numerical Linear Algebra}}
\urlstyle{same}
\setlength{\parindent}{0pt}
\setlength{\parskip}{5pt plus 1pt minus 1pt}
\linespread{1.04}
\setlength{\emergencystretch}{3em}
\widowpenalty=10000
\clubpenalty=10000
\displaywidowpenalty=10000
\allowdisplaybreaks[1]
\setlist{nosep,leftmargin=1.5em}
\providecommand{\tightlist}{\setlength{\itemsep}{0pt}\setlength{\parskip}{0pt}}
\providecommand{\pandocbounded}[1]{#1}
\pagestyle{fancy}
\fancyhf{}
\fancyhead[L]{\sffamily\footnotesize\color{muted}OPEN PROBLEMS IN NUMERICAL LINEAR ALGEBRA}
\fancyhead[R]{\sffamily\bfseries\footnotesize\color{accent}AV-02}
\fancyfoot[L]{\parbox[b]{0.9\textwidth}{\sffamily\fontsize{8}{10}\selectfont\color{muted}Editorial ratings; a bounded search cannot certify that a problem remains unsolved.\\Literature check: 2026-09-08}}
\fancyfoot[R]{\sffamily\footnotesize\color{muted}\thepage}
\renewcommand{\headrulewidth}{0.3pt}
\renewcommand{\headrule}{\hbox to\headwidth{\color{accent}\leaders\hrule height \headrulewidth\hfill}}
\makeatletter
\renewcommand{\section}{\Needspace{5\baselineskip}\@startsection{section}{1}{0pt}{12pt plus 2pt minus 2pt}{4pt}{\sffamily\bfseries\large\color{ink}}}
\renewcommand{\subsection}{\Needspace{4\baselineskip}\@startsection{subsection}{2}{0pt}{9pt plus 1pt minus 1pt}{3pt}{\sffamily\bfseries\normalsize\color{ink}}}
\makeatother
\setcounter{secnumdepth}{0}
\begin{document}
{\sffamily\small\color{muted}Intervals and absolute value equations\par}
\vspace{3pt}
{\raggedright\hyphenpenalty=10000\exhyphenpenalty=10000\sffamily\bfseries\fontsize{21}{24}\selectfont\color{ink}Hardness
of the spectral-norm condition number\par}
\vspace{7pt}
{\sffamily\small\textbf{Difficulty:} challenging\quad\textbf{Importance:} interesting
to the community\par}
\vspace{4pt}
{\color{accent}\hrule height 0.8pt}
\vspace{6pt}
\textbf{Status:} open; checked 2026-09-08\\
\textbf{Area:} condition estimation and rigorous error bounds

\subsection{Context and notation}\label{context-and-notation}

Absolute values and vector inequalities are componentwise. All
algorithmic questions use rational input in binary and the deterministic
Turing model. Input length includes the matrix dimensions and the bit
lengths of all rational coefficients.

For a real square matrix \(A\), define the diagonal perturbation family

\[
\mathcal D(A)=\{A-\operatorname{diag}(d):d\in[-1,1]^n\}.
\]

Call this family \textbf{regular} when every member is nonsingular. Only
diagonal entries vary. This is precisely the entrywise interval matrix
\([A-I_n,A+I_n]\).

\subsection{Problem statement}\label{problem-statement}

Let \(n\ge1\) and let \(A\in\mathbb Q^{n\times n}\) be promised to have
regular \(\mathcal D(A)\). Define

\[
c_2(A)=\max_{d\in[-1,1]^n}
\left\|(A-\operatorname{diag}(d))^{-1}\right\|_2
=\left(\min_{d\in[-1,1]^n}
\sigma_{\min}(A-\operatorname{diag}(d))\right)^{-1}.
\]

\subsubsection{Question}\label{question}

Is deciding \(c_2(A)\ge t\), given such an \(A\) and a positive rational
threshold \(t\), \(\mathsf{NP}\)-hard under polynomial-time Turing
reductions that query only inputs satisfying the regularity promise?

This is a decision formulation of the published hardness conjecture.
Equality belongs to the yes case. Checking the promise is outside the
task. The spectral norm is the operator norm induced by the Euclidean
vector norm.

\subsection{References}\label{references}

Moslem Zamani and Milan Hladík,
\href{https://doi.org/10.1007/s10107-021-01756-6}{\emph{Error bounds and
a condition number for the absolute value equations}}, Mathematical
Programming \textbf{198} (2023), 85--113, §2, paragraph before
Proposition 3; §2.1 and §7. The maximum is attained at a sign vector by
Proposition 2; Proposition 6 gives a formula for symmetric \(A\). The
conjecture concerns general matrices.

The quantity bounds forward error by residual for \(Ax-|x|=b\); see
Theorem 7. Its computation remains explicitly unclassified in Hladík et
al.'s \href{https://doi.org/10.1007/s10589-025-00717-5}{2026 survey},
§5.4.

\subsection{Status evidence}\label{status-evidence}

The original paper's latest arXiv version is
\href{https://arxiv.org/abs/1912.12904}{v2, January 17, 2020}; the
journal text was published January 30, 2022. Searches for ``absolute
value equations condition number 2026'', ``spectral norm complexity'',
and ``2-norm NP-hard'' found no resolution.
\end{document}
